\section{Limitations}\label{sec:limitations}

\emph{A ghost is a single copy.} Every chunk is stored once, in one repository, and GhostVault never
makes a second copy. Placement spreads storage; it does not protect against losing a repository, and
distributed placement makes a file depend on every repository that holds one of its chunks. What
GitHub deletes, loses or withholds cannot be brought back, and the user's use of GitHub is subject to
GitHub's terms, including its limits on repository size, request rates and bandwidth. Files that
cannot be lost need another copy elsewhere.

\emph{Removed is not erased.} Restoring a ghost removes its chunks from the branch, not from Git's
history, and GitHub may keep unreachable objects for a time. Every earlier version of the repository
backup stays in history too. A removed file stays unreadable only while its key exists nowhere. Getting
rid of old objects for good requires deleting the repository, or rewriting its history and asking
GitHub to purge cached objects. On SSDs and copy-on-write file systems, deleting the original does not
reliably destroy its old content either.

\emph{Keys at rest and in memory.} File keys are stored unencrypted in owner-only files
(\secref{sec:keys}), so anyone who can read the data folder, or a backup of it, and the repositories
can decrypt the files. Keys, passphrases and tokens pass through process memory, which Python cannot
reliably erase; swap, hibernation files and crash dumps may hold them.

\emph{The deletion window on Linux and macOS.} A program that already had the original open for
writing can still write to it in the moment between GhostVault's last check and the deletion
(\secref{sec:remove}). Windows closes that window; Linux and macOS offer no way to.

\emph{Metadata.} GitHub learns what \tabref{tab:learn} lists: how much is stored, approximate sizes,
and when ghosts come and go. A preview's existence reveals that a ghost is probably an image.

\emph{Throughput.} GitHub's pacing and limits, not GhostVault, bound how fast files can be ghosted
(\secref{sec:performance}); a large file takes minutes to hours, and GhostVault waits whenever GitHub
asks it to. Making a ghost also costs time linear in the number of ghosts.

\emph{One computer at a time.} A data folder is used by one GhostVault process at a time. Two
computers can share repositories, and the backup detects ghosts it does not know, but they do not
share or merge ghosts.

\emph{Scope.} Only repositories on \code{github.com} are supported, one file is ghosted at a time,
and folders are not. Tk exposes little to screen readers; every action can be reached from the
keyboard, but GhostVault cannot do more without leaving Tk. Standalone builds are unsigned unless a
signing certificate is configured, and there is no automatic update.

\emph{Verification.} The limits of what was verified are listed in \secref{sec:status}: the Windows
and macOS behaviour, the real-GitHub tests, the release workflow and the standalone applications were
not run for the version described here.
